\scanpage{12}
\deleted{puis schéma de groupes $G$, de type fini sur $\mathbb Z$, qui \deleted{\unsure{repr.}} est représenté matriciellement par} \unsure{$\mathcal G$|Rq.|$\mathbb E_{\mathbb Z}$|$\underline O_{\mathbb Z}$}. Considérons le sschéma en groupes de $\unsure{\mathrm{Gl}(3)|\mathrm{Gl}(n)}_{\mathbb Z}$ qui correspond aux matrices de la forme
\[
\begin{pmatrix}1&\alpha&\gamma\\0&\lambda&\beta\\0&0&\mu\end{pmatrix}
\]
(avec $\lambda\mu$ une unité bien sur), \inserted{il opère sur $M$} de la façon évidente qui consiste à interpréter $ABCDEF$ comme coeff. d'une forme
\[
\deleted{\unsure{X^2+AXY^2|X^3+AXY^2}} \quad X^2+AY^3+BXY^2+CY^2+DX+EY+F
\]
et en faisant sur $X,Y$, la substitution affine
\[
\deleted{X}Y'=\lambda\deleted{X}Y+\alpha,\quad X'=\mu X+\beta Y+\gamma
\]
[on peut aussi dire --- ce sera plus clair, que $G$ opère dans le vectoriel $\mathbb E_3$, donc dans $\Phi_3(\mathbb E_3)$, et qu'il laisse invariant le sous-\deleted{\unsure{espace}}\inserted{module loc. libre} de \deleted{\unsure{formes}} rang 6 des formes qui sont de la forme dite, d'où les opérations de $G$ sur $M$.

\margin{\uuline{Prop}}

Alors le foncteur qui à tt $S$ en-dessous associe \inserted{l'\unsure{ens|ens.} des classes à isom près} des courbes elliptiques ``structurées'' comme dessus sur $S$, est \uline{représentable} par $M$. De plus, \deleted{\unsure{p}}pour que deux élém de
% archive sheet number: 11
